\documentclass[11pt]{article}

\usepackage[a4paper,margin=1in]{geometry}

% ---- 中文支持：必须在 hyperref 之前加载 ----
\usepackage[fontset=macnew]{ctex}

\usepackage{amsmath,amssymb,amsfonts,bm,mathtools}
\usepackage{booktabs,multirow,array,tabularx,longtable,makecell}
\usepackage{placeins}

\usepackage[expansion=false]{microtype}

\usepackage{graphicx}
\PassOptionsToPackage{table}{xcolor}
\usepackage{xcolor}
\usepackage[table]{xcolor}
\usepackage{enumitem}
\usepackage{hyperref}
\usepackage[capitalize,noabbrev]{cleveref}
\usepackage[backend=bibtex,sorting=none,maxbibnames=99]{biblatex}
\DeclareFieldFormat{url}{{\small\url{#1}}}
\addbibresource{references.bib}
\usepackage{microtype}
\usepackage{authblk}
\usepackage{threeparttable}
\usepackage{tabularray}
\UseTblrLibrary{booktabs}
\usepackage{leftindex}

\usepackage{float}
\usepackage{tikz}
\usetikzlibrary{arrows.meta, calc, backgrounds, positioning}
\pgfdeclarelayer{foreground}
\pgfsetlayers{background, main, foreground}

\renewcommand{\topfraction}{0.92}
\renewcommand{\textfraction}{0.05}
\renewcommand{\floatpagefraction}{0.85}

\hypersetup{
    colorlinks=true,
    linkcolor=blue!50!black,
    citecolor=blue!50!black,
    urlcolor=blue!50!black
}

% ---- cleveref 中文化（必须在 cleveref 加载之后设置） ----
\crefname{section}{节}{节}
\crefname{subsection}{节}{节}
\crefname{subsubsection}{节}{节}
\crefname{figure}{图}{图}
\crefname{table}{表}{表}
\crefname{equation}{式}{式}
\crefname{algorithm}{算法}{算法}
\crefname{theorem}{定理}{定理}
\crefname{appendix}{附录}{附录}
\crefname{listing}{清单}{清单}
\crefname{definition}{定义}{定义}
\crefname{lemma}{引理}{引理}
\crefname{proposition}{命题}{命题}
\crefname{corollary}{推论}{推论}
\crefname{remark}{注记}{注记}
\crefname{assumption}{假设}{假设}
\Crefname{section}{节}{节}
\Crefname{subsection}{节}{节}
\Crefname{subsubsection}{节}{节}
\Crefname{figure}{图}{图}
\Crefname{table}{表}{表}
\Crefname{equation}{式}{式}
\Crefname{algorithm}{算法}{算法}
\Crefname{theorem}{定理}{定理}
\Crefname{appendix}{附录}{附录}
\Crefname{listing}{清单}{清单}
\Crefname{definition}{定义}{定义}
\Crefname{lemma}{引理}{引理}
\Crefname{proposition}{命题}{命题}
\Crefname{corollary}{推论}{推论}
\Crefname{remark}{注记}{注记}
\Crefname{assumption}{假设}{假设}

% ---- biblatex 参考文献标题中文化 ----
\DefineBibliographyStrings{english}{%
  bibliography = {参考文献},
  references   = {参考文献},
  andothers    = {等},
  page         = {页},
  pages        = {页},
}

% ---- 固定名称中文化（兜底，ctex 已设置大部分） ----
\renewcommand{\figurename}{图}
\renewcommand{\tablename}{表}
\renewcommand{\contentsname}{目录}
\renewcommand{\abstractname}{摘要}
\renewcommand{\refname}{参考文献}
\renewcommand{\appendixname}{附录}
% 注：article 类无 \bibname，不重定义
\AtBeginDocument{\renewcommand{\abstractname}{摘要}}

% Custom commands for email in authblk
\newcommand{\email}[1]{\texttt{#1}}
\definecolor{revisionmagenta}{RGB}{190,70,145}
\newcommand{\rev}[1]{\textcolor{revisionmagenta}{#1}}

\newcommand{\R}{\mathbb{R}}
\newcommand{\C}{\mathbb{C}}
\newcommand{\X}{\mathcal{X}}
\newcommand{\Y}{\mathcal{Y}}
\newcommand{\W}{\mathcal{W}}
\newcommand{\A}{\mathcal{A}}
\newcommand{\T}{\mathcal{T}}
\newcommand{\E}{\mathcal{E}}
\newcommand{\G}{\mathcal{G}}
\newcommand{\ten}[1]{\mathcal{#1}}
\newcommand{\mat}[1]{\mathbf{#1}}
\newcommand{\vecb}[1]{\mathbf{#1}}
\newcommand{\rank}{\operatorname{rank}}
\newcommand{\TTD}{\mathrm{TT}}
\newcommand{\cp}{\mathrm{CP}}
\newcommand{\bt}{\mathrm{BT}}
\newcommand{\mps}{\mathrm{MPS}}
\newcommand{\mha}{\mathrm{MHA}}
\newcommand{\ffn}{\mathrm{FFN}}
\newcommand{\kv}{\mathrm{KV}}
\newcommand{\lora}{\mathrm{LoRA}}
\newcommand{\peft}{\mathrm{PEFT}}
\newcommand{\set}[1]{\mathrm{#1}}
\newcommand{\llama}{\text{LLaMA-3-8B}}

\title{面向语言模型的张量方法：从词元表示到训练、\\ 适配、压缩、推理与可解释性}

%\author{}
\date{}

\author[1]{Matvei Tarasov}
 \affil[1]{Innopolis University, Innopolis 420500, Russia\\ \email{m.tarasov@innopolis.university}}

 \author[1]{Salman Ahmadi-Asl}
 \affil[1]{Innopolis University, Innopolis 420500, Russia\\ \email{s.ahmadiasl@innopolis.ru}}

 \author[2]{André L. F. de Almeida}
 \affil[2]{Department of Teleinformatics Engineering, Federal University of Cear\'a, Fortaleza 60455-760, Brazil\\ \email{andre@gtel.ufc.br}}

 \author[3]{Andrzej Cichocki}
 \affil[3]{Systems Research Institute of Polish Academy of Science and Warsaw University of Technology, Poland\\ \email{cichockiand@gmail.com}}

%\date{\today}
